\documentclass[../../main.tex]{subfiles}
\begin{document}

\begin{center}
  {\fontsize{48pt}{58pt}\selectfont \textbf{2013}\normalsize\textbf{年度}}
\end{center}

\begin{table}[htb]
  \centering
  \caption{出題テーマと難易度}
  \begin{tabular}{|c|p{6.2cm}|c|c|c|}
    \hline
    \textbf{問題番号} & \textbf{テーマ} & \textbf{思考量} & \textbf{計算量} & \textbf{総合} \\
    \hline
    第1問 & 複素数で表す平面上の点の周期性 & 普 & 易 & 易 \\
    \hline
    第2問 & 二つの関数のグラフの共有点の個数 & 易 & 易 & 易 \\
    \hline
    第3問 & コイン投げによる得点と勝敗の確率 & 普 & 普 & 普 \\
    \hline
    第4問 & 三角形の内部の点と単位ベクトルの和 & 易 & 普 & 普 \\
    \hline
    第5問 & 連続する整数の積の十進表示 & 難 & 易 & 難 \\
    \hline
    第6問 & 正方形の回転でできる立体の共通部分 & 普 & 普 & 普 \\
    \hline
  \end{tabular}
\end{table}

\begin{table}[htb]
  \centering
  \caption{解答}
  \renewcommand{\arraystretch}{1.4}
  \begin{tabular}{|c|c|p{9cm}|}
    \hline
    \textbf{問題番号} & \textbf{小問} & \textbf{答え} \\
    \hline
    第1問 & -- & $(a,b)=\left(\dfrac{1}{2},
    \pm\dfrac{\sqrt{3}}{2}\right)$ \\
    \hline
    第2問 & -- & $a=-\dfrac{2}{5\pi}$，または
    $\dfrac{2}{5\pi}<a<\dfrac{2}{3\pi}$ \\
    \hline
    第3問 & (1) & $n=2m$なら$p(n)=\dfrac{m}{2^{2m}}$，
    $n$が奇数なら$p(n)=\dfrac{(n-1)(n-3)}{2^{n+2}}$ \\
    \cline{2-3}
    & (2) & $\displaystyle\sum_{n=1}^{\infty}p(n)=\dfrac{16}{27}$ \\
    \hline
    第4問 & (1) & $\angle APB=\angle APC=\dfrac{2\pi}{3}$ \\
    \cline{2-3}
    & (2) & $|PA|=\dfrac{\sqrt{7}}{7}$，
    $|PB|=\dfrac{2\sqrt{7}}{7}$，
    $|PC|=\dfrac{4\sqrt{7}}{7}$ \\
    \hline
    第5問 & (1) & $x>\dfrac{1}{2}\left(3y^2-1+
    \sqrt{(3y^2-1)^2+4y(y^2-1)}\right)$ \\
    \cline{2-3}
    & (2) & 条件を満たす連続3整数の積が存在する（証明問題） \\
    \hline
    第6問 & (1) & 切り口の面積は
    $S(t)=\dfrac{8\sqrt{2}}{3}\sqrt{1-t^2}$ \\
    \cline{2-3}
    & (2) & 共通部分の体積は$\dfrac{32\sqrt{2}}{9}$ \\
    \hline
  \end{tabular}
\end{table}

\paragraph{試験形態}
6問．試験時間は要確認．

\end{document}
